% Esquisse d'un programme (cote 119, copy 2): archive cote 119, site page 100 (PDF page 101), typescript page 22.
% typescript. Draft compared: esquisse/pages/esq_119_024.tex (copy 1, site 24).
% AI-transcribed by Claude Opus 5.5 (claude-opus-5-5, adaptive thinking, effort high), Message Batches API, batch esq119_full (msgbatch_015xqDHkTumDysRF1qngZhRJ), 2026-11-13; not yet checked by eye.
\scanpage{119/100}
\mspage{22}

polyèdre régulier envisagé, est fini. Dans le cas du corps de
base $\mathbb{R}$ (ou $\mathbb{C}$, ce qui revient au même), et pour $n = 2$, les cas
finis sont bien connus depuis l'antiquité - ce qui n'exclut pas
que le point de vue schématique y fasse apparaître des charmes
nouveau ; on peut dire cependant qu'en spécialisant l'isosaèdre
(par exemple) sur des corps de base finis de caractéristique
arbitraire, c'est toujours un icosaèdre, avec sa combinatoire
propre et le même groupe d'automorphismes simple d'ordre 60 qu'on
obtient. La même remarque s'applique aux polyèdres réguliers finis
de dimension supérieure, étudiés de façon systématique dans deux
beaux livres de Coxeter. La situation est toute autre si on part
d'un polyèdre régulier \uline{infini}, sur un corps tel que $\mathbb{Q}$ disons,
et qu'on le ``spécialise'' sur les corps premiers $\mathbb{F}_p$ (opération
bien définie pour tout $p$ sauf un nombre fini de nombres premiers).
Il est clair que tout polyèdre régulier sur un corps fini est fini -
\uline{on trouve donc une infinité de polyèdres réguliers finis pour} $p$
\uline{variable}, \uline{dont le type combinatoire}, ou ce qui revient au même,
le groupe des automorphismes, \uline{varie de façon} ``\uline{arithmétique}'' avec p.
Cette situation est particulièrement intrigante  dans le cas où
$n = 2$, où on dispose de la relation explicitée au paragraphe pré-%
cédent entre 2-cartes combinatoires, et courbes algébriques dé-%
finies sur des corps de nombres. Dans ce cas, un polyèdre régu-%
lier infini défini sur un corps infini quelconque (et de ce fait,
sur une sous-$\mathbb{Z}$-algèbre à deux générateurs de celui-ci) donne donc
naissance à une infinité de courbes algébriques définies sur des
corps de nombres, qui sont des revêtements galoisiens ramifiés
seulement en 0, 1, $\infty$ de la droite projective standard. Le cas opti-%
mum est bien sûr celui où on part du 2-polyèdre régulier universel,
ou plutôt de celui qui s'en déduit par passage au corps des frac-%
tions $\mathbb{Q}(\alpha_1,\alpha_2)$ de son anneau de base. Ceci soulève une foule de
questions nouvelles, aussi bien des vagues que des précises, dont
je n'ai eu loisir encore d'examiner de plus près aucune - je ne
citerai que celle-ci : quelles sont exactement les 2-cartes régu-%
lières finies, ou ce qui revient au même, les groupes quotients

\margin{100}
